\clearpage
\subsection{Flag Meanings}
Integer condition codes are held in FLAGS. Integer instructions leave FLAGS unchanged unless their instruction
description explicitly defines a FLAGS write. Every FLAGS-writing instruction replaces the complete (:ref:base.registers.STATE.FLAGS.Z:), (:ref:base.registers.STATE.FLAGS.N:), (:ref:base.registers.STATE.FLAGS.C:), and (:ref:base.registers.STATE.FLAGS.V:)
image; partial FLAGS writes do not exist. An instruction that does not write FLAGS preserves the complete image.

For an operand width of n bits, (:term:base.terms.ordinary:) integer ALU results are evaluated modulo 2\textasciicircum{}n unless the
instruction explicitly defines a wider intermediate.

Explicit FLAGS-writing integer instructions include CMP, TEST, ADC, SBB, INCF, DECF, BTEST, BSET, BCLR, BCHG,
SETF, WRFLAGS, CMPXCHG, SEGLEA, and the bounds-check instructions. CMPJcc and TESTJcc compute temporary condition flags
for their branch decision and do not update architectural FLAGS. ADD, SUB, logic operations, INC, DEC, shifts,
rotates, moves, extensions, min/max, and multiply/divide instructions leave FLAGS unchanged.

\BedrockTableCaption{Integer Condition-Code Bits}
\begin{BedrockTabular}{@{}>{\raggedright\arraybackslash}p{0.45in}>{\raggedright\arraybackslash}p{1.05in}>{\raggedright\arraybackslash}p{3.90in}@{}}
\toprule
\rowcolor{BedrockHeaderFill}
\textbf{Bit} & \textbf{Name} & \textbf{Meaning}\\
\midrule
(:ref:base.registers.STATE.FLAGS.Z:) & zero & Set when the result value is zero.\\
(:ref:base.registers.STATE.FLAGS.N:) & negative & Set from the most significant result bit at the operand size.\\
(:ref:base.registers.STATE.FLAGS.C:) & carry/borrow & Set on unsigned carry out for addition or unsigned borrow for subtraction.\\
(:ref:base.registers.STATE.FLAGS.V:) & overflow & Set on signed overflow or an explicitly reported exceptional condition.\\
\bottomrule
\end{BedrockTabular}
